\scanpage{89}
% archive sheet number: 88
% block text | line 1
Or on voit que
% block display | line 2
\[
h(\unsure{\gamma|\sigma})=g(\gamma)f(\gamma)^{-1}
\]
% block text | lines 3-8
est un $1$-cocycle de $\Gamma$ à valeurs dans $A_f$. On est ramené à prouver qu'il est \deleted{trivial}\inserted{$\sim0$}, sachant qu'il est \deleted{trivial}\inserted{$\sim0$} sur $\Gamma_i$. \unsure{Aut.} nous revenons donc au cas où l'un des cocycles $f,g$ est l'unité, soit $g$. Mais dans ce cas, on utilise la suite exacte
% block diagram | line 9
\diagram{\begin{tikzcd}[ampersand replacement=\&,column sep=small]
e\arrow[r,dash]\&H^1(\Gamma/\Gamma_i,G^{\Gamma_i})\arrow[r,dash]\&H^1(\Gamma,G)\arrow[r,dash]\&H^1(\Gamma_i,G)
\end{tikzcd}}
% block text | line 10
et on est ramené à prouver que
% block display | line 11
\[
H^1(\Gamma/\Gamma_i,G^{\Gamma_i})=\{e\}
\]
% block text | lines 12-20
Or $G^{\Gamma_i}=(A^{\Gamma_i})^*$. Soit $\Gamma'=\Gamma/\Gamma_i$, $A'=A^{\Gamma_i}$, \deleted{donc} $O'=O^{\Gamma_i}$. On sait que $O'$ est \inserted{local} noethérien, que son corps de restes est isomorphe à celui de $O$ donc que $\Gamma/\Gamma_i=\Gamma'$ y opère \uline{fidèlement}. De plus on sait que $O$ est un $O'$-module noethérien, d'où résulte qu'il en est de $=$ de $A$, donc de $A'$. $A'$ est une $O'$-algèbre, et $\Gamma'$ y opère de façon compatible avec les opérations sur $O'$. On est ramené à prouver
% block display | line 21
\[
H^1(\Gamma,A^*)=(e)
\]
% block text | line 22
\deleted{\unsure{si $O$}} dans le cas particulier où $\Gamma_i=\{e\}$,
